%%%%%%%%%%%%%%%%%%%%%%%%%%%%%%%%%%%%%%%%%%%%%%%%%%%%%%%%%%%%%%%%%%%%%%%%%%%%%%%%%%%%%%%%%%
%%	
%% File Name:		dhbw_theme_example.tex				
%%											
%% Description:		Example presentation using the dhbw beamer theme.
%%
%% Author:			Hannes Bartle																				
%% 					DHBW Ravensburg Campus Friedrichshafen		
%%					September 2016											
%%
%%%%%%%%%%%%%%%%%%%%%%%%%%%%%%%%%%%%%%%%%%%%%%%%%%%%%%%%%%%%%%%%%%%%%%%%%%%%%%%%%%%%%%%%%%
\documentclass{beamer}


\usepackage[utf8]{inputenc}
\usepackage[T1]{fontenc}


\title{There Is No Largest Prime Number}
\date[ISPN ’80]{27th International Symposium of Prime Numbers}
\author[Euclid]{Euclid of Alexandria \texttt{euclid@alexandria.edu}}
\institute{DHBW Ravensburg Campus Friedrichshafen}

\usetheme{dhbw}
\renewcommand{\usecompanylogo}{1}
%\setlength{\companylogox}{2cm}
%\setlength{\companylogoy}{0.9\paperheight}
\newcommand{\printcompanylogo}{\includegraphics[width=4cm]{example/airbus_logo.png}}




\begin{document}

	\begin{frame}
		\titlepage
	\end{frame}


	\begin{frame} 

		\frametitle{There Is No Largest Prime Number} 
		\framesubtitle{The proof uses \textit{reductio ad absurdum}.} 

		\begin{theorem}
			There is no largest prime number. 
		\end{theorem} 
		\begin{enumerate} 
			\item<1-> Suppose $p$ were the largest prime number. 
			\item<2-> Let $q$ be the product of the first $p$ numbers. 
			\item<3-> Then $q+1$ is not divisible by any of them. 
			\item<1-> But $q + 1$ is greater than $1$, thus divisible by some prime number not in the first $p$ numbers.
		\end{enumerate}
	\end{frame}

	\begin{frame}{A longer title}
		\begin{itemize}
			\item one
			\item two
		\end{itemize}
	\end{frame}

\end{document}